\section{Method}
\label{sec:method}

Given a pair of images and camera calibration, our method reconstructs scale-aware point maps with an affinity learning in the first stage.
In the second stage, we project such point maps onto the target view and refine this geometry to anchor Gaussian primitives for the rendering task. 
An overview of our 2D-coarse-to-3D-fine pipeline is shown in Fig.~\ref{fig:pipeline}. 

\subsection{Scale-Aware Geometry Reconstruction}
\label{sec:geometry}
\paragraph{Point Map.} It is introduced by DUSt3R~\cite{wang2024dust3r} as a novel but scale-invariant representation $X \in [0, 1]^{W \times H \times 3}$ of 3D scene, which is associated with corresponding image $I$ of resolution $W \times H$. 
We apply the follow-up, MASt3R~\cite{leroy2024grounding}, to predict two pieces of point map from source views, $i \in \{l, r\}$ for left and right view, to represent the scene.
As a coarse stage, we take $W = 512$ and $H=288$.
MASt3R encodes the input images into patches of features $F^{i}$ with ViT and then decodes them into $X^i$ in canonical space.
Without any stereo constraint, it is capable of predicting reasonable geometries from two cameras in a large sparsity, but hard to control the target view with real camera parameters, leading to jitters when inferring consecutive frames.
%
In the following, we transform point maps from canonical space to real space in an absolute scale with an affinity, in the format of scaling $S \in \mathbb{R}^3$ and translation $T \in \mathbb{R}^{W \times H \times 3}$.

\paragraph{Scaling.} It is a global factor related to camera intrinsic parameters such as focal $f$. In addition, the distance $d$ between two cameras provides a cue of measurement in real space. Thus, we embed them with positional encoding~\cite{mildenhall2020nerf} $\text{PE}$
\begin{equation}
    e = \text{PE}(f, d)
\end{equation}
We further process the encoded features $F$ with self- $Att_s$ and cross-attention $Att_c$
\begin{align}
\begin{split}
\label{formula:att}
    \langle\mathbf{Q,K,V}\rangle &= \langle F \mathbf{W}^Q, F \mathbf{W}^K, F \mathbf{W}^V \rangle \\
    f_s &= \text{Avg}(Att_s(\mathbf{Q}_l, \mathbf{K}_l, \mathbf{V}_l)) \\
    f_c &= \text{Avg}(Att_c(\mathbf{Q}_l, \mathbf{K}_r, \mathbf{V}_r))
\end{split}
\end{align}
where the average operator is used to extract global information. We use an $MLP$ to compute the scaling factor
\begin{equation}
    S = MLP(f_s, f_c, e)
    \label{eq:scale}
\end{equation}
We note that the degree of freedom of $S$ is 3 to deal with the distortion of original point map reconstruction of MASt3R.

\paragraph{Translation.} Although the point maps could be rescaled to real space by multiplying the scaling factor, there could still exist point-wise shifts due to the lack of stereo constraint in MASt3R. 
Inspired by the view consistency check~\cite{yan2020dense} in MVS, we believe that such a shift depends not only on the features in one view but also on the corresponding features in the other view. 
Thus, we process the aforementioned ViT features with a lightweight convolutional encoder $\mathcal{E}_c$ to yield feature map $F_c^i = \mathcal{E}(F^i)$. Furthermore, we obtain the feature map $F^{j \rightarrow i}_c$ in the other view $j$, $j \in \{ l, r \}$ and $i \neq j$, by first projecting the rescaled points $SX^i$ onto the view $j$ and then querying the corresponding features with bilinear sampling
\begin{equation}
    F^{j \rightarrow i}_c = \text{Query}(F^j_c , \text{Proj}(SX^i, K^j))
\end{equation}
where $K$ is the camera parameter.
We follow the idea of iterative updating~\cite{teed2020raft,lipson2021raft-stereo} to compute point-wise translation with $\text{GRU}$ operator~\cite{cho2014properties} by considering feature maps from both views and the position of each point
\begin{equation}
    T^i = \text{GRU}(F^i, F^{j \rightarrow i}, SX^i)
\end{equation}
We obtain the position of the point set in real space with the learned affine transform
\begin{equation}
    X_{t}^i = SX^i + T^i
\end{equation}

\subsection{Rendering via Gaussian Plane}
\label{sec:rendering}

\paragraph{3D Refinement.} When splatting the aforementioned point set to the target view, we still observe inevitable jitters and holes due to the lack of 3D stereo constraint, see Fig.~\ref{fig:abla_full}(a). 
%
Given the paired images $I_f$ in fine resolution of $W = 1024$ and $H=576$, we encode them with convolutional layers $\mathcal{E}_f$ into $F_f^i = \mathcal{E}_f(I_f^i)$. 
%
Additionally, we project the transformed point set $X_t$ onto the target view $k$ with $\alpha$-blending mechanism in Gaussian Splatting to yield the initial depth map $\mathcal{D}^k$.
%
For each pixel $(u,v)$, we sample several position candidates $\{d_1, d_2,...d_N\}$ near the initial depth value $d = \mathcal{D}(u,v)$ along camera ray.
%
For each candidate $d(u,v,n)$, we warp the feature from the source views to the target view
\begin{equation}
    p^k(u,v,n) = \text{Proj}^{-1}(d^k(u,v,n), K^k)
\end{equation}
\begin{equation}
\label{eq:warp}
    F_f^{i \rightarrow k}(u,v,n) = \text{Query}(F_f^i , \text{Proj}(p^k(u,v,n), K^i))
\end{equation}
where the warping process can be efficiently achieved by matrix operation.
We further process the aggregation of features with 3D convolutions $\mathcal{E}_{3D}$ into a feature volume 
\begin{equation}
    \Phi^k = \mathcal{E}_{3D}(F_f^{l \rightarrow k} , F_f^{r \rightarrow k})
\end{equation}
Following ENeRF~\cite{lin2022enerf}, we compute the depth probability distribution $w_n$ along the camera ray by regressing with the feature volume $\Phi$.
%
The final position of Gaussian primitives can be represented with the refined depth $\bar{d} = \Sigma_n w_n d_n$.

\paragraph{Gaussian Plane.} Once the position of Gaussian is determined, Gaussian plane $\mathcal{G}$ consists of four attribute maps of color, rotation, scaling and opacity
\begin{equation}
    \mathcal{G} = \{ \mathcal{P}_c, \mathcal{P}_r, \mathcal{P}_s, \mathcal{P}_o \}
\end{equation}
Using the warping process in Eq.~\ref{eq:warp}, we obtain color $\{ C^{l \rightarrow k}, C^{r \rightarrow k} \}$ of target view warped from source views.
We further query the feature $\phi$ from feature volume $\Phi$ for Gaussian primitives via tri-linear interpolation.
%
We thus learn a weighted color to initialize the Gaussian color
\begin{align}
    w^{i}_{c} &= MLP_{c}(f_f^l, f_f^r, \phi)\\
    C^k &= \Sigma_{i} w_c^i C^{i \rightarrow k} 
    \label{eq:color_stage2}
\end{align}
We arrange all features into the format of a 2D map and further aggregate them into the feature map $\mathcal{M} = \text{Agg}\{ f_f^l, f_f^r, \phi \}$.
Following GPS-Gaussian, we yield rotation, scaling, and opacity map via convolutional heads $h_a, a=\{r, s, o \}$, considering the Gaussian position $Y$
\begin{equation}
    \mathcal{P}_a = h_a (\mathcal{M}, Y)
\label{eq:ha}
\end{equation}
In addition, we update the initial color with a learned residual color map 
\begin{align}
    \Delta C &= {h}_c(\mathcal{M}, Y, C)\\
    \mathcal{P}_c &= \alpha C + (1-\alpha) \Delta C
    \label{eq:color}
\end{align}
Finally, we splat the Gaussian plane $\mathcal{G}$ in a fine resolution of $1024 \times 576$ to render an image $\hat{I}$ in a higher resolution of $1280 \times 720$.


\begin{table*}[ht!]
\small

\centering
\begin{tabular}{l|ccc|ccc|ccc}
\toprule
\multicolumn{1}{c}{\multirow{1}[6]{*}{Method}} & \multicolumn{3}{c}{Camera} & \multicolumn{3}{c}{GoPro} & \multicolumn{3}{c}{Mobile} \\

\cmidrule{2-10} \multicolumn{1}{c}{} & PSNR$\uparrow$ & SSIM$\uparrow$ & \multicolumn{1}{c}{LPIPS$\downarrow$} & PSNR$\uparrow$ & SSIM$\uparrow$  & \multicolumn{1}{c}{LPIPS$\downarrow$} & PSNR$\uparrow$ & SSIM$\uparrow$ & \multicolumn{1}{c}{LPIPS$\downarrow$} \\ \midrule

NoPoSplat  & 25.035 & 0.866 & 0.173 & 26.128 & 0.889 & 0.121 & 21.594 & 0.591 & \underline{0.272}  \\
4D-GS  & 27.814 & 0.906 & 0.150 & 27.244 & 0.907 & 0.205 & 25.655 & \underline{0.825} & 0.284  \\
MVSplat     & 27.899 & 0.902 & 0.148 & \underline{29.942} & 0.934 & 0.157 & \textbf{26.545} & 0.805 &  0.314  \\
MVSGaussian & \underline{29.326} & \underline{0.957} & \textbf{0.069} & 27.413 & 0.926 & 0.151 & 19.927 & 0.683 & \underline{0.272}   \\
ENeRF  & 28.272 & 0.943 & 0.084 & 29.906 & \underline{0.943} & \underline{0.108} & 20.579 & 0.640 & 0.302   \\
Ours   & \textbf{32.220} & \textbf{0.957} & \underline{0.079} & \textbf{31.640} & \textbf{0.955} & \textbf{0.096} & \underline{25.721} & \textbf{0.827} & \textbf{0.244} \\
\bottomrule
\end{tabular}
\vspace{-1mm}
\caption{\textbf{Quantitative comparison of rendering methods on multi-view datasets.} NoPoSplat~\cite{ye2024nopose}, MVSplat~\cite{chen2024mvsplat} and MVSGaussian~\cite{liu2024mvsgaussian} are feed-forward Gaussian Splatting methods and ENeRF~\cite{lin2022enerf} is feed-forward NeRF based method, while 4D-GS~\cite{Wu2024_4dgs} is optimization based 4D Gaussian Splatting method. \textbf{Bold} highlights the top-performing method, while \underline{underline} indicates suboptimal performance across various evaluation criteria.}
\vspace{-2mm}
\label{tab:num_compare_bg}
\end{table*}


\begin{figure*}[ht!]
  \centering
  \includegraphics[width=0.88\textwidth]{figs/4_comp.pdf}
  \vspace{-2mm}
  \caption{\textbf{Quantitative comparison of rendering.} We show results of (a) ENeRF~\cite{lin2022enerf}, (b) MVSGaussian~\cite{liu2024mvsgaussian}, (c) MVSplat~\cite{chen2024mvsplat}, (d) Ours and (e) Ground Truth for GoPro, Camera and Mobile datasets.}
  \vspace{-4mm}
  \label{fig:quantative_comp}
\end{figure*}

\subsection{Training}

We define the rendering loss as a combination of L1 loss $\mathcal{L}_1$ and SSIM loss~\cite{wang2004ssim} $\mathcal{L}_{ssim}$
\begin{equation}
    \mathcal{L}_{render}(\hat{I}, I^{gt}) = \beta_1 \mathcal{L}_1 + \beta_2 \mathcal{L}_{ssim}
    \label{eq:render_loss}
\end{equation}
where $\hat{I}$ and $I^{gt}$ stand for rendering image and ground truth image.
Since the transformed geometry largely impacts the rendering module, we propose a 2-stage training strategy. 


\paragraph{Stage 1.} We firstly train the affine transformed point maps with captured multi-view images in a self-supervised manner without any 3D supervision.
%
To this end, we predict Gaussian planes $\mathcal{G}^i = \{ \mathcal{P}_p^i, \mathcal{P}_c^i, \mathcal{P}_r^i, \mathcal{P}_s^i, \mathcal{P}_o^i \}, i=\{ l,r \}$ on source views. Among them, the position and color plane can be obtained with point map $X^i_t$ and input image $I^i$. 
%
We further use the auxiliary layers $\bar{h}_a$, similar to convolutional operator $h_a$ in Eq.~\ref{eq:ha}, to predict $\mathcal{P}_a^i, a=\{r, s, o\}$.
%
Once the auxiliary Gaussian planes are done, we can render the image on the target view.  
Inspired by GPS-Gaussian+~\cite{zhou2024gps}, we propose a regularization term as Chamfer distance between two 6-dimensional point sets $P^l, P^r$
\vspace{-2mm}
\begin{align}
\begin{split}
    CD(i \rightarrow j) &= \frac{1}{|{P}^i|}\sum_{p_i \in {P}^i} \min_{p_j \in {P}^j} \|p_i - p_j\|_2 \\ 
    \mathcal{L}_{CD} &= CD(l \rightarrow r) + CD(r \rightarrow l)
\end{split}
\label{eq:geo_regular}
\end{align}
where $p_i$ is pixel-wise point on point map $X_t^i(u,v)$ associated with the corresponding pixel color $I^i(u,v)$.
Such a regularization term allows two pieces of point maps to converge to a better geometry.

Therefore, we supervise the affinity learning and auxiliary layers with rendering loss and the regularization term
\begin{equation}
    \mathcal{L}_{stage1} = \mathcal{L}_{render} + \gamma\mathcal{L}_{CD}
    \label{eq:loss1}
\end{equation}
Note that during the training, we freeze the weight of the MASt3R~\cite{leroy2024grounding} network and no longer require the geometry ground truth. 

\paragraph{Stage 2.} The scale-aware point maps in the previous step allow us to initialize the depth of the target view, which largely improves the stability of the training process in stage 2. Specifically, we train the 3D refinement module and Gaussian planes with photo-metric loss. In practice, we have an initial color plane $\hat{I}_f = C$ in fine resolution and a splatting image $\hat{I}_h$ in a higher resolution.
So the training loss is the combination of two rendering losses 
\begin{equation}
    \mathcal{L}_{stage2} = \lambda_1 \mathcal{L}_{render}(\hat{I}_f, I_{f}^{gt}) + \lambda_2 \mathcal{L}_{render}(\hat{I}_h, I_{h}^{gt})
    \label{eq:loss2}
\end{equation}

For both stages 1 and 2, we do not require 3D geometry supervision, which facilitates the training process on real captured 2D images.